\section{Introduction}
\label{sec:introduction}

On a smooth Riemannian manifold, a lower bound for sectional curvature
is equivalent to a lower comparison bound for geodesic triangles.
Gigli's distributional curvature tensor makes the analytic side of
this statement meaningful on spaces with synthetic lower Ricci
curvature bounds. The purpose of this paper is to prove that the same
equivalence holds in that setting, with the Hausdorff measure and
dimension hypotheses proposed by Gigli.



An Alexandrov lower curvature bound \(\CBB(\kappa)\) means that
geodesic triangles satisfy lower comparison with the simply connected
surface of constant curvature \(\kappa\), at the usual model scale
when \(\kappa>0\). A \(\RCD(K,N)\) space combines the
unreduced optimal-transport curvature-dimension condition
\(\CD(K,N)\) with quadratic Cheeger energy. The latter gives a
Hilbertian Sobolev calculus in which the tensor below is defined.
We take the reference measure to be the exact Hausdorff measure,
not merely a measure equivalent to it.

On a full-support finite-dimensional \(\RCD\) space \((M,d,m)\),
we use the nonpositive Laplacian:
\[
\int_M\langle\nabla f,\nabla h\rangle\dd m
=-\int_M(\Delta f)h\dd m.
\]
The domain \(D(\Delta)\) is the domain of the \(L^2\) generator.
In particular, its members belong to \(W^{1,2}(M)\). Our test classes are
\begin{align}
\label{eq:test-class}
\Test(M)
&=\left\{f\in D(\Delta):
\begin{array}{l}f\text{ has a bounded globally Lipschitz representative},\\
\Delta f\in W^{1,2}(M)
\end{array}\right\},\\
\label{eq:test-vector-class}
\TestV(M)
&=\left\{\sum_{j=1}^{\ell}a_j\nabla b_j:
\ell<\infty,\ a_j,b_j\in\Test(M)\right\}.
\end{align}
For test fields set
\([X,Y]=\nabla_XY-\nabla_YX\), where \(\nabla\) is the weak
Levi--Civita covariant derivative. The curvature functional is
\begin{align}
\label{eq:curvature-definition}
R(X,Y,Z,W)(f)=\int_M&
-\langle\nabla_X(fW),\nabla_YZ\rangle
-f\langle\nabla_YZ,W\rangle\divg X \notag\\
&+\langle\nabla_Y(fW),\nabla_XZ\rangle
+f\langle\nabla_XZ,W\rangle\divg Y \notag\\
&-f\langle\nabla_{[X,Y]}Z,W\rangle\dd m.
\end{align}
The last derivative means contraction of the weak derivative of \(Z\)
with \([X,Y]\). This is the sign convention corresponding to
\(\nabla_X\nabla_Y-\nabla_Y\nabla_X-\nabla_{[X,Y]}\).
All terms in \eqref{eq:curvature-definition} are integrable in the
test classes \eqref{eq:test-class}--\eqref{eq:test-vector-class}.
No classical curvature tensor is assumed.

\begin{theorem}[Gigli's characterization]
\label{main:equivalence}
Let \(n\ge2\) be an integer, let \(\kappa\in\R\), and let \((M,d)\)
be a complete separable metric space. The following conditions are
equivalent.
\begin{enumerate}[label=\textup{(\Alph*)}]
\item \((M,d)\) is an \(n\)-dimensional Alexandrov space with curvature
bounded below by \(\kappa\).
\item With \(m=\mathcal H^n_d\), the space \((M,d,m)\) has full support,
is \(\RCD((n-1)\kappa,n)\), and satisfies
\begin{equation}
\label{eq:main-curvature}
R(X,Y,Y,X)(f)\ge
\kappa\int_M f\bigl(|X|^2|Y|^2-\langle X,Y\rangle^2\bigr)\dd m
\end{equation}
for every \(X,Y\in\TestV(M)\) and every nonnegative \(f\in\Test(M)\).
\end{enumerate}
The \(\RCD\) condition is unreduced, and the test classes and
curvature sign are exactly those defined above.
\end{theorem}

\subsection{History and significance}

Lott--Villani and Sturm independently developed synthetic lower
Ricci curvature bounds through convexity properties of entropy along
Wasserstein geodesics \cite{LottVillani2009,Sturm2006I,Sturm2006II}.
Ambrosio--Gigli--Savar\'e introduced the \(\RCD(K,\infty)\)
framework, imposing quadratic Cheeger energy to select a Riemannian
class \cite{AGS2014}. Ambrosio--Gigli--Mondino--Rajala extended
the theory to reference measures finite on bounded sets and removed
the strengthened entropy condition in the original definition
\cite{AGMR2015}. These developments supply the ambient Ricci-curvature
and Sobolev framework for the sectional-curvature question considered here.

Petrunin and Zhang--Zhu established the connection between
Alexandrov lower curvature bounds and the curvature-dimension condition
\cite{Petrunin2011,ZhangZhu2010}.
The quadratic Sobolev theory of Kuwae--Machigashira--Shioya
\cite{KMS2001} supplies the Riemannian structure.
The full-scope Alexandrov-to-\(\RCD\) statement used below is recorded
in \cite[Theorem~2.9 and Corollary~2.10]{KapovitchKettererSturm}.
Gigli subsequently defined
\eqref{eq:curvature-definition} and conjectured the
characterization proved here in dimensions at least two
\cite[Conjecture~1.1]{Gigli2019}.
His construction gives tensor symmetries and tensoriality, but does
not by itself relate a lower tensor bound to triangle comparison.

Theorem~\ref{main:equivalence} resolves this conjecture positively
for every integer \(n\ge2\).
The exact choice of reference measure is important:
the assumption \(m=\mathcal H^n\) supplies noncollapsed structure,
including \(n\)-dimensional Euclidean tangents almost everywhere
\cite{DPG2018}.
The use of the original test classes is important as well.
The theorem is neither a statement about a classical tensor on a
regular subset nor a restriction to compactly supported fields.

Optimal transport also gives a characterization of sectional
curvature in the smooth setting. Ketterer--Mondino express sectional
and intermediate Ricci bounds through entropy inequalities along
Wasserstein geodesics supported on lower-dimensional rectifiable
sets \cite{KettererMondino2018}. The transport calculation here pairs
full-dimensional densities with Gigli's distributional tensor.

Several related developments locate the analytic difficulties.
Lebedeva--Petrunin construct a weak curvature tensor for smoothable
Alexandrov spaces, using noncollapsing smooth approximations with a
common sectional lower bound \cite{LebedevaPetrunin2024}. Their weak-limit
tensor is a separate framework from the global test-class functional
considered here.
Brena--Gigli studied measure representations of weak Hessians and
curvature tensors, and explained the regularity obstruction in
passing from a distributional Hessian bound to convexity
\cite[Remark~2.4 and Remark~2.18]{BrenaGigli2024}.
For \(C^1\) Riemannian metrics, Er\"os--Kunzinger--Ohanyan--Vardabasso
proved a local distributional-to-Alexandrov comparison result
\cite[Theorem~3.1]{EKOV2026}.
The argument here works directly with the Sobolev calculus of a
noncollapsed \(\RCD\) space.

\subsection{The new arguments}

Section~\ref{sec:forward} proves the forward implication.
The converse proceeds through the index inequality of
Theorem~\ref{heat:index}, the Hessian comparison of
Theorem~\ref{frames:hessian}, and the passage to comparison along
every geodesic in Section~\ref{sec:weighted-comparison}.

\paragraph{Integrated configurations.}
Section~\ref{sec:forward} extracts the curvature functional from
small Alexandrov configurations. The expansion is integrated against
transported densities. A quantitative first-order estimate for the
interpolating density permits pairing with fixed \(L^2\) Hessian
coefficients. The geometric comparison and the transport expansion
bound the same energy difference. Averaging cancels the lower-order
terms. Suitable endpoint corrections then identify the remaining
coefficient as the curvature term plus a nonnegative square error.
Local choices of gradient tests make this error arbitrarily small.
Tensoriality, approximation and exhaustion then
recover the global test classes in Theorem~\ref{fwd:curvature}.
The calculation uses only the weak derivatives provided by the test calculus.

\paragraph{A heat-regularized index inequality.}
Section~\ref{sec:heat-index} pairs a forward heat factor with a backward
heat factor normalized by the prescribed terminal density.
Their product solves a continuity equation. The logarithmic
potentials yield a Riccati identity, in which the terms containing
the inverse viscosity cancel. The remaining acceleration error must
be controlled uniformly up to the terminal time, including the
singular measure terms in the Bochner inequality. A one-sided
estimate supplies this control. An action identity gives strong convergence of
the velocities in the density-weighted norm, which is needed for
the quadratic index energy.

\paragraph{Frames through exceptional times.}
Almost-everywhere spatial regularity alone does not give a frame at
every time on a path. Section~\ref{sec:frames} upgrades harmonic
splitting maps to frames along almost every relevant entire path. A Hausdorff-content
estimate controls the regular points where a countable family of
test gradients fails to have full rank. Capacity estimates and
bounded-compression path estimates remove the exceptional sets.
Deng's continuity theorem for tangent cones along every geodesic
\cite{Deng2020} supplies regularity at all interior times.
Caputo--Gigli--Pasqualetto established parallel transport along test
plans in the noncollapsed setting \cite{CGP2022}. Here finite-dimensional
Sobolev ODEs in the selected frames construct and approximate the
Jacobi trials needed by the index inequality. Only a countable trial catalog
is needed; an integrated density theorem for parallel fields is not
assumed.

\paragraph{From Hessian comparison to every geodesic.}
Section~\ref{sec:weighted-comparison} constructs bounded global
modifications of the distance function. Their distributional Hessian
bounds fit the general theorem of the companion paper
\cite[Theorem~1.1]{WeakHessianCompanion}, recalled with its exact
hypotheses in Theorem~\ref{weight:every-geodesic}. That theorem applies
to every full-support \(\RCD(K,N)\) space with \(1<N<\infty\),
including collapsed spaces. Its proof uses bounded changes of measure
and a closed weighted Bochner inequality; it does not use the
characterization or the frame construction of this paper. Applied to
the modified distances, it gives the model differential inequality
along every minimizing geodesic. A one-dimensional comparison and
the local-to-global Alexandrov theorem finish the argument.

Figure~\ref{fig:proof-route} summarizes the reverse implication and
the point where the companion theorem is used.
\begin{figure}[t]
\centering
\begin{tikzpicture}[
 box/.style={draw=blue!40!black,rounded corners=2pt,align=center,
 text width=39mm,minimum height=15mm,font=\small,inner sep=3pt},
 >={Stealth[length=2mm]}]
\node[box] (tensor) at (0,0) {Global distributional\\sectional bound};
\node[box] (index) at (4.8,0) {Heat index inequality\\on fixed terminal paths};
\node[box] (hessian) at (9.6,0) {Radial Hessian bound\\as a signed measure};
\node[box] (distance) at (9.6,-2.1) {Bounded global\\distance modifications};
\node[box] (ode) at (4.8,-2.1) {Model ODE along\\every geodesic};
\node[box] (cbb) at (0,-2.1) {Triangle comparison\\and global $\CBB(\kappa)$};
\draw[->] (tensor) -- (index);
\draw[->] (index) -- (hessian);
\draw[->] (hessian) -- (distance);
\draw[->] (distance) -- node[below=9mm,font=\scriptsize,align=center]
 {Theorem~\ref{weight:every-geodesic}} (ode);
\draw[->] (ode) -- (cbb);
\end{tikzpicture}
\caption{The reverse implication. The noncollapsed frame and countable
trial construction connects the heat index estimate to the radial
Hessian bound. The companion theorem supplies the passage to every
geodesic after the bounded distance modification and has general
finite-dimensional RCD scope.}
\label{fig:proof-route}
\end{figure}


The frame and pathwise trial constructions in the first two
subsections of Section~\ref{sec:frames} also have an independent scope.
They apply on full-support noncollapsed \(\RCD(K,n)\) spaces,
with \(m=\mathcal H^n\) and integer \(n\ge2\), to optimal geodesic
plans with bounded positions, bounded speed and bounded compression
throughout the relevant interval. The approximation is pathwise,
using a countable catalog of global trials.

We state the foundational \(\RCD\) results in the form used in the
proof and give complete arguments for the additional transport,
index, frame and distance-comparison steps. The general reweighting
argument is supplied by the companion paper \cite{WeakHessianCompanion}.
Numbered external results refer to the source versions linked in
the bibliography; version notes specify differences from published
numbering where necessary.
